% \documentclass[border=15pt]{standalone}
% \usepackage{tikz,amsmath}
% \usetikzlibrary{arrows.meta,decorations.pathreplacing}

% % ========== 配色方案 ==========
% \definecolor{hidFill}{RGB}{199,203,243}
% \definecolor{hidLine}{RGB}{52,58,143}
% \definecolor{inFill}{RGB}{183,224,177}
% \definecolor{inLine}{RGB}{32,102,32}
% \definecolor{outFill}{RGB}{243,199,199}
% \definecolor{outLine}{RGB}{143,52,52}

% \begin{document}
% \begin{tikzpicture}[
%     font=\small, line join=round,
%     every node/.style={inner sep=0pt, outer sep=0pt, align=center}
%   ]

%   % ========== 通用 3D 堆叠切片宏 ==========
%   % #1: x, #2: y, #3: 网格格数, #4: 单元格尺寸, #5: 颜色前缀, #6: 叠加层数, #7: 层间深度步长
%   \newcommand{\stackbox}[7]{%
%     \foreach \k in {#6,...,1}{%
%       \pgfmathtruncatemacro{\op}{round(100 - (55 * \k / #6))}%
%       \fill[#5Fill!\op, draw=#5Line, line width=.35pt]
%             ({#1+#7*\k},{#2+#7*\k}) rectangle ++(#3*#4,#3*#4);
%     }
%     \fill[#5Fill] (#1,#2) rectangle ++(#3*#4,#3*#4);
%     \draw[step=#4, #5Line!60, line width=.35pt, shift={(#1,#2)}] (0,0) grid (#3*#4,#3*#4);
%     \draw[#5Line, thick] (#1,#2) rectangle ++(#3*#4,#3*#4);
%   }

%   % ========== 红色 3D 感受野/卷积立体框宏 ==========
%   % #1: x, #2: y, #3: 宽, #4: 高, #5: 深度偏移
%   \newcommand{\prism}[5]{%
%     \draw[outLine, dashed, line width=.7pt]
%           (#1+#5,#2+#5) rectangle (#1+#5+#3,#2+#5+#4);
%     \draw[outLine, line width=1.1pt]
%           (#1+#5,#2+#5+#4) -- (#1+#5+#3,#2+#5+#4);
%     \draw[outLine, dashed, line width=.7pt]
%           (#1,#2)--(#1+#5,#2+#5) (#1+#3,#2)--(#1+#3+#5,#2+#5);
%     \draw[outLine, line width=1.1pt] (#1,#2) rectangle (#1+#3,#2+#4);
%     \draw[outLine, line width=1.1pt]
%           (#1,#2+#4)--(#1+#5,#2+#5+#4) (#1+#3,#2+#4)--(#1+#3+#5,#2+#5+#4);
%   }

%   % 1. 左侧：总输入特征图 (7x7x24)
%   \stackbox{-5.4}{-1.225}{7}{0.35}{in}{10}{0.08}
%   \prism{-5.4}{0.175}{1.05}{1.05}{0.8}
%   \node[align=center] at (-4.0, -2.1) {\textbf{Input Feature Map}\\$H\times W\times C_{in}$\\($7\times 7\times 24$)};

%   % 加长后的 Split 箭头与说明文字 (指向左侧大花括号尖端)
%   \draw[->, >=Stealth, line width=1.2pt, inLine] (-2.0, 0) -- (-0.5, 0);
%   \node[above, font=\bfseries\footnotesize, inLine] at (-1.4, 0.18) {Split\\($g=3$)};

%   % 2. 左侧大花括号（开口向右，紧凑包裹三组分支，无多余箭头）
%   \draw[inLine, line width=1.2pt, decorate, decoration={brace, amplitude=12pt}]
%        (-0.05, -5.3) -- (-0.05, 5.3);

%   % 3. 中间三组卷积分支
%   \foreach \Y/\grp in {4.6/1, 0.0/2, -4.6/3}{
%     % 分组标题
%     \node[font=\bfseries\small, hidLine] at (1.6, \Y + 1.9) {Group \grp};

%     % (1) 分组输入切片 7x7
%     \stackbox{0.4}{\Y - 1.225}{7}{0.35}{in}{5}{0.09}
%     \prism{0.4}{\Y + 0.175}{1.05}{1.05}{0.45}
%     \node[font=\scriptsize] at (1.6, \Y - 1.6) {$7\times 7\times 8$};

%     % (2) 卷积符号 *
%     \node[font=\LARGE] at (3.85, \Y) {$\ast$};

%     % (3) 紫色卷积核组 (每组画2个核，中间为 \cdots 与倍数)
%     % 核 1 (上)
%     \stackbox{4.4}{\Y + 0.45}{3}{0.28}{hid}{5}{0.09}
%     \prism{4.4}{\Y + 0.45}{0.84}{0.84}{0.45}

%     % 中间省略号与倍数
%     \node[font=\large] at (4.82, \Y + 0.2) {$\cdots$};
%     \node[font=\scriptsize, anchor=west] at (5.45, \Y - 1.2) {$\times \frac{N}{3}$};

%     % 核 2 (下)
%     \stackbox{4.4}{\Y - 1.29}{3}{0.28}{hid}{5}{0.09}
%     \node[font=\scriptsize] at (4.82, \Y - 1.65) {$3\times 3\times 8$};

%     % (4) 等号 =
%     \node[font=\LARGE] at (6.7, \Y) {$=$};

%     % (5) 分组输出切片 5x5 (已去除指向右花括号的箭头)
%     \stackbox{7.3}{\Y - 0.875}{5}{0.35}{out}{5}{0.09}
%     % 左上角高亮输出单元
%     \fill[outLine!40] (7.3, \Y + 0.525) rectangle ++(0.35, 0.35);
%     \draw[outLine, line width=1pt] (7.3, \Y + 0.525) rectangle ++(0.35, 0.35);
%     \node[font=\scriptsize] at (8.15, \Y - 1.3) {$5\times 5\times \frac{N}{3}$};
%   }

%   % 4. 右侧大花括号（开口向左，紧贴三组输出，无多余箭头）
%   \draw[outLine, line width=1.2pt, decorate, decoration={brace, amplitude=12pt}]
%        (9.9, 5.3) -- (9.9, -5.3);

%   % 从右花括号尖端指向 Concat 模块的箭头
%   \draw[->, >=Stealth, line width=1.2pt, outLine] (10.4, 0) -- (11.0, 0);

%   % Concat 处理节点
%   \node[draw=outLine, fill=outFill!30, rounded corners=3pt, line width=1pt,
%         inner sep=5pt, font=\bfseries\small, text=outLine] (concat) at (11.8, 0) {Concat};

%   % 从 Concat 指向最终总输出图的箭头
%   \draw[->, >=Stealth, line width=1.2pt, outLine] (concat.east) -- (13.3, 0);

%   % 5. 最右侧：合并后的最终输出特征图 (5x5x24)
%   \stackbox{13.6}{-0.875}{5}{0.35}{out}{10}{0.08}
%   % 对应的合成高亮输出像素点
%   \fill[outLine!40] (13.6, 0.525) rectangle ++(0.35, 0.35);
%   \draw[outLine, line width=1pt] (13.6, 0.525) rectangle ++(0.35, 0.35);
%   \node[align=center] at (14.5, -2.1) {\textbf{Output Feature Map}\\$H'\times W'\times C_{out}$\\($5\times 5\times N$)};

% \end{tikzpicture}
% \end{document}
\documentclass[border=15pt]{standalone}
\usepackage{tikz,amsmath}
\usetikzlibrary{arrows.meta,decorations.pathreplacing}

% ========== 配色方案 ==========
\definecolor{hidFill}{RGB}{199,203,243}
\definecolor{hidLine}{RGB}{52,58,143}
\definecolor{inFill}{RGB}{183,224,177}       % 真实特征图绿色
\definecolor{inLine}{RGB}{32,102,32}
\definecolor{outFill}{RGB}{243,199,199}
\definecolor{outLine}{RGB}{143,52,52}
\definecolor{padFill}{RGB}{228,244,226}     % 浅浅绿色填充 (Padding)
\definecolor{padLine}{RGB}{120,170,120}     % padding 网格与虚线边框色 (柔和草绿)

\begin{document}
\begin{tikzpicture}[
    font=\small, line join=round,
    every node/.style={inner sep=0pt, outer sep=0pt, align=center}
  ]

  % ========== 通用 3D 堆叠切片宏 (无 Padding) ==========
  % #1: x, #2: y, #3: 网格格数, #4: 单元格尺寸, #5: 颜色前缀, #6: 叠加层数, #7: 层间深度步长
  \newcommand{\stackbox}[7]{%
    \foreach \k in {#6,...,1}{%
      \pgfmathtruncatemacro{\op}{round(100 - (55 * \k / #6))}%
      \fill[#5Fill!\op, draw=#5Line, line width=.35pt]
            ({#1+#7*\k},{#2+#7*\k}) rectangle ++(#3*#4,#3*#4);
    }
    \fill[#5Fill] (#1,#2) rectangle ++(#3*#4,#3*#4);
    \draw[step=#4, #5Line!60, line width=.35pt, shift={(#1,#2)}] (0,0) grid (#3*#4,#3*#4);
    \draw[#5Line, thick] (#1,#2) rectangle ++(#3*#4,#3*#4);
  }

  % ========== 带 Padding 的 3D 堆叠切片宏 ==========
  % #1: x, #2: y, #3: 内层真实数据格数, #4: 每格尺寸, #5: 颜色前缀, #6: 叠加层数, #7: 深度步长, #8: padding圈数
  \newcommand{\stackboxpad}[8]{%
    \pgfmathtruncatemacro{\totgrid}{#3 + 2*#8}
    \pgfmathsetmacro{\totdim}{\totgrid * #4}
    \pgfmathsetmacro{\indim}{#3 * #4}
    \pgfmathsetmacro{\padoff}{#8 * #4}
    % --- 背后叠加层 (从后往前绘制) ---
    \foreach \k in {#6,...,1}{%
      \pgfmathtruncatemacro{\op}{round(100 - (55 * \k / #6))}%
      % 背后 padding 浅浅绿色底色
      \fill[padFill!\op, draw=padLine, line width=.35pt]
            ({#1+#7*\k},{#2+#7*\k}) rectangle ++(\totdim,\totdim);
      % 背后 真实特征图区域
      \fill[#5Fill!\op, draw=#5Line, line width=.35pt]
            ({#1+#7*\k+\padoff},{#2+#7*\k+\padoff}) rectangle ++(\indim,\indim);
    }
    % --- 正面层 ---
    % 1. 正面 padding 底色与虚线网格
    \fill[padFill] (#1,#2) rectangle ++(\totdim,\totdim);
    \draw[step=#4, padLine!70, dashed, line width=.35pt, shift={(#1,#2)}] (0,0) grid (\totdim,\totdim);
    % 2. 正面 真实数据特征图底色与实线网格
    \fill[#5Fill] ({#1+\padoff},{#2+\padoff}) rectangle ++(\indim,\indim);
    \draw[step=#4, #5Line!60, line width=.35pt, shift={({#1+\padoff},{#2+\padoff})}] (0,0) grid (\indim,\indim);
    % 3. 边框：内层特征图实线框，外圈 padding 虚线框
    \draw[#5Line, thick] ({#1+\padoff},{#2+\padoff}) rectangle ++(\indim,\indim);
    \draw[padLine, thick, dashed] (#1,#2) rectangle ++(\totdim,\totdim);
    % 4. 正面 padding 单元格填入 0
    \pgfmathtruncatemacro{\lastidx}{\totgrid-1}
    \foreach \ix in {0,...,\lastidx}{%
      \node[font=\fontsize{5}{6}\selectfont, text=padLine!90] at ({#1+(\ix+0.5)*#4}, {#2+0.5*#4}) {0};
      \node[font=\fontsize{5}{6}\selectfont, text=padLine!90] at ({#1+(\ix+0.5)*#4}, {#2+(\lastidx+0.5)*#4}) {0};
    }
    \pgfmathtruncatemacro{\midend}{\lastidx-1}
    \foreach \iy in {1,...,\midend}{%
      \node[font=\fontsize{5}{6}\selectfont, text=padLine!90] at ({#1+0.5*#4}, {#2+(\iy+0.5)*#4}) {0};
      \node[font=\fontsize{5}{6}\selectfont, text=padLine!90] at ({#1+(\lastidx+0.5)*#4}, {#2+(\iy+0.5)*#4}) {0};
    }
  }

  % ========== 红色 3D 感受野/卷积立体框宏 ==========
  % #1: x, #2: y, #3: 宽, #4: 高, #5: 深度偏移
  \newcommand{\prism}[5]{%
    \draw[outLine, dashed, line width=.7pt]
          (#1+#5,#2+#5) rectangle (#1+#5+#3,#2+#5+#4);
    \draw[outLine, line width=1.1pt]
          (#1+#5,#2+#5+#4) -- (#1+#5+#3,#2+#5+#4);
    \draw[outLine, dashed, line width=.7pt]
          (#1,#2)--(#1+#5,#2+#5) (#1+#3,#2)--(#1+#3+#5,#2+#5);
    \draw[outLine, line width=1.1pt] (#1,#2) rectangle (#1+#3,#2+#4);
    \draw[outLine, line width=1.1pt]
          (#1,#2+#4)--(#1+#5,#2+#5+#4) (#1+#3,#2+#4)--(#1+#3+#5,#2+#5+#4);
  }

  % 1. 左侧：总输入特征图 (内层 5x5 + Padding 1 = 7x7x24)
  \stackboxpad{-5.4}{-1.225}{5}{0.35}{in}{10}{0.08}{1}
  \prism{-5.4}{0.175}{1.05}{1.05}{0.8}
  \node[align=center] at (-4.0, -2.15) {\textbf{Input Feature Map}\\$H\times W\times C_{in}$\\($5\times 5\times 24$)};

  % 加长后的 Split 箭头与说明文字 (指向左侧大花括号尖端)
  \draw[->, >=Stealth, line width=1.2pt, inLine] (-2.0, 0) -- (-0.5, 0);
  \node[above, font=\bfseries\footnotesize, inLine] at (-1.4, 0.18) {Split\\($g=3$)};

  % 2. 左侧大花括号（开口向右，紧凑包裹三组分支）
  \draw[inLine, line width=1.2pt, decorate, decoration={brace, amplitude=12pt}]
       (-0.05, -5.3) -- (-0.05, 5.3);

  % 3. 中间三组卷积分支
  \foreach \Y/\grp in {4.6/1, 0.0/2, -4.6/3}{
    % 分组标题
    \node[font=\bfseries\small, hidLine] at (1.6, \Y + 1.9) {Group \grp};

    % (1) 分组输入切片 (内层 5x5 + Padding 1 = 7x7x8)
    \stackboxpad{0.4}{\Y - 1.225}{5}{0.35}{in}{5}{0.09}{1}
    \prism{0.4}{\Y + 0.175}{1.05}{1.05}{0.45}
    \node[font=\scriptsize] at (1.6, \Y - 1.65) {$5\times 5\times 8$};

    % (2) 卷积符号 *
    \node[font=\LARGE] at (3.85, \Y) {$\ast$};

    % (3) 紫色卷积核组 (每组画2个核，中间为 \cdots 与倍数)
    % 核 1 (上)
    \stackbox{4.4}{\Y + 0.45}{3}{0.28}{hid}{5}{0.09}
    \prism{4.4}{\Y + 0.45}{0.84}{0.84}{0.45}

    % 中间省略号与倍数
    \node[font=\large] at (4.82, \Y + 0.2) {$\cdots$};
    \node[font=\scriptsize, anchor=west] at (5.45, \Y - 1.2) {$\times \frac{N}{3}$};

    % 核 2 (下)
    \stackbox{4.4}{\Y - 1.29}{3}{0.28}{hid}{5}{0.09}
    \node[font=\scriptsize] at (4.82, \Y - 1.65) {$3\times 3\times 8$};

    % (4) 等号 =
    \node[font=\LARGE] at (6.7, \Y) {$=$};

    % (5) 分组输出切片 5x5
    \stackbox{7.3}{\Y - 0.875}{5}{0.35}{out}{5}{0.09}
    % 左上角高亮输出单元
    \fill[outLine!40] (7.3, \Y + 0.525) rectangle ++(0.35, 0.35);
    \draw[outLine, line width=1pt] (7.3, \Y + 0.525) rectangle ++(0.35, 0.35);
    \node[font=\scriptsize] at (8.15, \Y - 1.3) {$5\times 5\times \frac{N}{3}$};
  }

  % 4. 右侧大花括号（开口向左，紧贴三组输出）
  \draw[outLine, line width=1.2pt, decorate, decoration={brace, amplitude=12pt}]
       (9.9, 5.3) -- (9.9, -5.3);

  % 从右花括号尖端指向 Concat 模块的箭头
  \draw[->, >=Stealth, line width=1.2pt, outLine] (10.4, 0) -- (11.0, 0);

  % Concat 处理节点
  \node[draw=outLine, fill=outFill!30, rounded corners=3pt, line width=1pt,
        inner sep=5pt, font=\bfseries\small, text=outLine] (concat) at (11.8, 0) {Concat};

  % 从 Concat 指向最终总输出图的箭头
  \draw[->, >=Stealth, line width=1.2pt, outLine] (concat.east) -- (13.3, 0);

  % 5. 最右侧：合并后的最终输出特征图 (5x5xN)
  \stackbox{13.6}{-0.875}{5}{0.35}{out}{10}{0.08}
  % 对应的合成高亮输出像素点
  \fill[outLine!40] (13.6, 0.525) rectangle ++(0.35, 0.35);
  \draw[outLine, line width=1pt] (13.6, 0.525) rectangle ++(0.35, 0.35);
  \node[align=center] at (14.5, -2.1) {\textbf{Output Feature Map}\\$H\times W\times C_{out}$\\($5\times 5\times N$)};

\end{tikzpicture}
\end{document}